\chapter{Arquitetura}
\label{cap:arquitetura}

O GamP-SCPeriop é um monólito: um único programa ASP.NET Core serve a aplicação cliente e a \gls{api}, e guarda tudo numa única base de dados. Este capítulo justifica essa escolha, descreve os componentes, explica como são resolvidas as rotas do cliente e da \gls{api}, e apresenta o catálogo completo de endpoints.

\section{Estilo arquitetural}
\label{sec:estilo}

Escolheu-se um monólito com três projetos, em vez de serviços independentes, pelas seguintes razões:

\begin{itemize}
    \item \textbf{Escala.} A aplicação serve uma instituição, com dezenas de supervisores e algumas centenas de estudantes. Não há partes do sistema com cargas tão diferentes que justifiquem escalá-las em separado.
    \item \textbf{Consistência.} As operações principais tocam muitas tabelas de uma só vez: criar um percurso a partir de um molde copia módulos, componentes e badges, e guardar uma avaliação atualiza o progresso e atribui badges. Numa única base de dados, estas operações são transações simples.
    \item \textbf{Equipa e operação.} Um único programa, desenvolvido por uma equipa pequena, é mais simples de testar, de pôr em produção e de manter.
    \item \textbf{Código partilhado.} Cliente e servidor usam as mesmas classes de entidades e de DTOs, definidas uma única vez na biblioteca partilhada.
\end{itemize}

O custo desta escolha é que todas as partes são publicadas em conjunto e que uma falha no servidor afeta toda a aplicação. Para o contexto em causa, esse custo é aceitável.

\section{Componentes}

A solução tem três projetos (Tabela~\ref{tab:projetos}). O servidor aloja o cliente: no primeiro acesso, o navegador descarrega a aplicação Blazor WebAssembly como ficheiros estáticos; a partir daí, a aplicação corre no navegador e só comunica com o servidor através da \gls{api} \gls{rest} (Figura~\ref{fig:arquitetura}).

\begin{table}[H]
    \centering
    \small
    \begin{tabularx}{\linewidth}{lX}
        \toprule
        \textbf{Projeto} & \textbf{Responsabilidade} \\
        \midrule
        \texttt{GamP-SCPeriop} & Cliente Blazor WebAssembly: páginas, layouts, componentes reutilizáveis, estado de autenticação e chamadas à \gls{api}. \\
        \midrule
        \texttt{GamP-SCPeriop.Server} & Servidor ASP.NET Core: aloja o cliente, expõe a \gls{api} \gls{rest}, valida os \gls{jwt}, acede à base de dados com Entity Framework Core e gera os relatórios \gls{pdf} com QuestPDF. \\
        \midrule
        \texttt{GamP-SCPeriop.Shared} & Biblioteca partilhada: entidades, DTOs e enumerações usados pelos dois lados. \\
        \bottomrule
    \end{tabularx}
    \caption{Projetos da solução.}
    \label{tab:projetos}
\end{table}

\begin{figure}[H]
    \centering
    \resizebox{\linewidth}{!}{% Arquitetura do monólito: cliente Blazor WebAssembly, servidor ASP.NET Core e dados.
\begin{tikzpicture}[
    comp/.style={draw=gInk, fill=white, rounded corners=3pt, font=\scriptsize, align=center,
      text width=3.9cm, minimum height=1.0cm, inner sep=3pt, line width=0.6pt},
    compsis/.style={comp, fill=gBrand!12},
    titulo/.style={font=\scriptsize\bfseries, text=gInk, anchor=west},
    seta/.style={->, >={Stealth[length=2mm]}, draw=gInk, line width=0.6pt},
    nota/.style={font=\tiny\itshape, text=gInkSoft, align=center},
  ]
  % ---------- Navegador ----------
  \node[titulo] (tc) at (-2.1,4.9) {Navegador (Blazor WebAssembly)};
  \node[comp] (c1) at (0,3.9)   {Páginas, layouts e componentes\\{\tiny\ttfamily Pages/ Layout/ Components/}};
  \node[comp] (c2) at (0,2.45)  {Router +\\AuthorizeRouteView\\{\tiny\ttfamily App.razor}};
  \node[comp] (c3) at (0,1.0)   {CustomAuthStateProvider\\{\tiny lê o perfil a partir do JWT}};
  \node[comp] (c4) at (0,-0.45) {HttpClient \emph{API} +\\AuthTokenHandler\\{\tiny junta o cabeçalho Bearer}};
  \node[comp, dashed] (c5) at (0,-1.9) {Session storage do navegador\\{\tiny guarda o JWT}};

  % ---------- Servidor ----------
  \node[titulo] (ts) at (5.4,4.9) {Servidor (ASP.NET Core, .NET 10)};
  \node[compsis] (s1) at (7.5,3.9)   {Ficheiros estáticos do cliente\\{\tiny\ttfamily index.html, .wasm, .dll}};
  \node[compsis] (s2) at (7.5,2.45)  {Autenticação JWT + autorização\\{\tiny\ttfamily [Authorize], [Authorize(Roles)]}};
  \node[compsis] (s3) at (7.5,1.0)   {14 controladores REST\\{\tiny\ttfamily /api/<controlador>/\ldots}};
  \node[compsis] (s4) at (7.5,-0.45) {Serviços\\{\tiny BadgeService,}\\{\tiny EvaluationPdfGenerator}};
  \node[compsis] (s5) at (7.5,-1.9)  {AppDbContext (EF Core)\\{\tiny 15 tabelas; migrações no arranque}};

  % ---------- Dados ----------
  \node[titulo] (td) at (11.75,4.9) {Dados};
  \node[comp, text width=2.6cm] (d2) at (13.3,1.0) {Pasta \texttt{SecureUploads/}\\{\tiny PDF das componentes}};
  \node[comp, text width=2.6cm, minimum height=1.3cm, cylinder, shape border rotate=90,
        aspect=0.18, rounded corners=0pt] (d1) at (13.3,-1.9) {SQL Server\\{\tiny LocalDB em desenvolvimento}};

  % ---------- Biblioteca partilhada ----------
  \node[comp, text width=9.4cm, fill=gCard] (sh) at (3.75,-3.75)
      {\textbf{GamP-SCPeriop.Shared}: entidades, DTOs e enumerações usados pelos dois lados};

  % ---------- Ligações ----------
  \draw[seta, dashed] (s1.west) -- node[nota, above] {1.º carregamento} (c1.east);
  \draw[seta] (c4.east) -- ++(1.55,0) |- (s2.west);
  \node[nota, anchor=south] at ($(s2.west)+(-0.9,0)$) {HTTP + JWT};
  \draw[seta] (s5.east) -- node[nota, above] {SQL} (d1.west);
  \draw[seta] (s3.east) -- node[nota, above] {ficheiros} (d2.west);
  \draw[seta] (c1) -- (c2);
  \draw[seta] (c2) -- (c3);
  \draw[seta] (c1.west) -- ++(-0.3,0) |- (c4.west);
  \draw[seta] (c4) -- (c5);
  \draw[seta] (s2) -- (s3);
  \draw[seta] (s3) -- (s4);
  \draw[seta] (s4) -- (s5);
  \draw[seta, dashed, gInkSoft] (sh.north -| c5.south) -- (c5.south);
  \draw[seta, dashed, gInkSoft] (sh.north -| s5.south) -- (s5.south);

  % ---------- Fundos ----------
  \begin{scope}[on background layer]
    \node[area=gInkSoft, fit=(tc)(c1)(c5), inner sep=6pt] {};
    \node[area=gBrandDark, fit=(ts)(s1)(s5), inner sep=6pt] {};
    \node[area=gPorIniciar, fit=(td)(d1)(d2), inner sep=6pt] {};
  \end{scope}
\end{tikzpicture}
}
    \caption{Arquitetura do GamP-SCPeriop. A azul, os componentes do servidor; a seta tracejada é o descarregamento inicial da aplicação.}
    \label{fig:arquitetura}
\end{figure}

\subsection{Pipeline do servidor}

Cada pedido atravessa, por esta ordem, os componentes configurados em \texttt{Program.cs}:

\begin{enumerate}
    \item \textbf{CORS}, que aceita pedidos de qualquer origem;
    \item \textbf{Redireção para HTTPS};
    \item \textbf{Ficheiros estáticos}, incluindo os ficheiros da aplicação Blazor (\texttt{UseBlazor\-Framework\-Files} e \texttt{Use\-Static\-Files});
    \item \textbf{Encaminhamento}, que escolhe o controlador e a ação;
    \item \textbf{Autenticação e autorização}, que validam o \gls{jwt} e os atributos \texttt{[Authorize]};
    \item \textbf{Controladores}, e por fim o \textbf{fallback}, que devolve \texttt{index.html} a qualquer pedido que não corresponda a um ficheiro nem a um controlador.
\end{enumerate}

No arranque, o servidor aplica as migrações pendentes à base de dados (32 até à data) e, se não existir nenhum utilizador, cria a conta de administrador inicial. Em desenvolvimento, publica também a documentação OpenAPI em \texttt{/swagger}.

\subsection{Serviços e lógica de negócio}

A lógica de negócio está repartida entre os controladores e dois serviços:

\begin{itemize}
    \item \texttt{BadgeService}: verifica e atribui as badges de módulo e de percurso, e cria as notificações correspondentes;
    \item \texttt{EvaluationPdfGenerator}: gera o relatório global de avaliação de um estudante.
\end{itemize}

As cópias de moldes e de percursos e o cálculo do progresso estão escritos diretamente nos controladores (\texttt{PathwayController}, \texttt{EnrollmentController} e \texttt{EvaluationController}). Mover esta lógica para serviços próprios tornaria o código mais fácil de testar, e é a primeira melhoria proposta para a arquitetura.

\section{Rotas}
\label{sec:rotas}

A aplicação tem dois conjuntos de rotas, resolvidos em lugares diferentes:

\begin{itemize}
    \item \textbf{Rotas do cliente}, como \texttt{/dashboard} ou \texttt{/pathway/7}. São declaradas em cada página com a diretiva \texttt{@page} e resolvidas no navegador pelo \emph{Router} do Blazor. O Capítulo~\ref{cap:sitemap} lista-as todas.
    \item \textbf{Rotas da \gls{api}}, todas começadas por \texttt{/api/}. São declaradas com atributos nos controladores e resolvidas no servidor.
\end{itemize}

A Figura~\ref{fig:rotas} mostra o percurso de um pedido. O servidor tenta, por ordem, servir um ficheiro estático e encontrar um controlador; se nenhum corresponder, devolve \texttt{index.html}, e a aplicação Blazor, já no navegador, escolhe a página.

\begin{figure}[H]
    \centering
    \resizebox{0.82\linewidth}{!}{% Como um URL é resolvido: primeiro pelo servidor, depois (se for o caso) pelo Router do Blazor.
\begin{tikzpicture}[
    passo/.style={draw=gInk, fill=white, rounded corners=3pt, font=\scriptsize, align=center,
      text width=3.0cm, minimum height=0.9cm, inner sep=2pt, line width=0.6pt},
    passosis/.style={passo, fill=gBrand!12},
    decisao/.style={draw=gInk, fill=white, diamond, aspect=2.2, font=\scriptsize, align=center,
      inner sep=1pt, text width=2.0cm, line width=0.6pt},
    fim/.style={passo, fill=gApprovedBg, draw=gApproved},
    erro/.style={passo, fill=gReprovedBg, draw=gReproved},
    seta/.style={->, >={Stealth[length=2mm]}, draw=gInk, line width=0.6pt},
    rot/.style={font=\tiny\itshape, text=gInkSoft, inner sep=1.5pt},
    zona/.style={font=\scriptsize\bfseries, text=gInk, anchor=west},
  ]
  \node[passo] (p0) at (0,0) {Pedido do navegador\\{\tiny\ttfamily GET /pathway/7}};

  % ---- Servidor ----
  \node[zona] (zs) at (1.0,-0.95) {No servidor (\texttt{Program.cs})};
  \node[decisao] (d1) at (0,-2.1) {Ficheiro estático existe?};
  \node[fim] (r1) at (4.4,-2.1) {Devolve o ficheiro\\{\tiny\ttfamily .js, .css, .wasm, .dll}};
  \node[decisao] (d2) at (0,-4.0) {Rota de um controlador?};
  \node[passosis] (r2) at (4.4,-4.0) {Controlador e ação\\{\tiny\ttfamily api/[controller] + [HttpGet(\ldots)]}};
  \node[passosis] (r2b) at (8.6,-4.0) {Valida o JWT e o\\atributo \texttt{[Authorize]}};
  \node[fim] (r2c) at (8.6,-2.1) {Resposta JSON\\{\tiny 200, 400, 401, 404}};
  \node[passosis] (r3) at (0,-5.8) {\texttt{MapFallbackToFile}\\devolve \texttt{index.html}};

  % ---- Cliente ----
  \node[zona] (zc) at (1.0,-7.45) {No navegador (\texttt{App.razor})};
  \node[decisao] (d3) at (0,-8.6) {Alguma página tem esta rota?};
  \node[erro] (r4) at (4.4,-8.6) {\emph{NotFound}: \\«Página não encontrada»};
  \node[decisao] (d4) at (0,-10.5) {Tem o perfil exigido?};
  \node[fim] (r5) at (4.4,-10.5) {Mostra a página\\{\tiny\ttfamily @page /pathway/\{Id:int\}}};
  \node[erro] (r6) at (0,-12.2) {\emph{NotAuthorized}:\\redireciona para \texttt{/}};

  \draw[seta] (p0) -- (d1);
  \draw[seta] (d1) -- node[rot, above] {sim} (r1);
  \draw[seta] (d1) -- node[rot, right] {não} (d2);
  \draw[seta] (d2) -- node[rot, above] {sim} (r2);
  \draw[seta] (r2) -- (r2b);
  \draw[seta] (r2b) -- (r2c);
  \draw[seta] (d2) -- node[rot, right] {não} (r3);
  \draw[seta] (r3) -- (d3);
  \draw[seta] (d3) -- node[rot, above] {não} (r4);
  \draw[seta] (d3) -- node[rot, right] {sim} (d4);
  \draw[seta] (d4) -- node[rot, above] {sim} (r5);
  \draw[seta] (d4) -- node[rot, right] {não} (r6);

  \begin{scope}[on background layer]
    \node[area=gBrandDark, fit=(zs)(d1)(r1)(r2b)(r3), inner sep=6pt] {};
    \node[area=gApproved, fit=(zc)(d3)(r4)(r5)(r6), inner sep=6pt] {};
  \end{scope}
\end{tikzpicture}
}
    \caption{Resolução de um URL: primeiro no servidor, depois no navegador.}
    \label{fig:rotas}
\end{figure}

\subsection{Rotas do cliente}

Cada página declara a sua rota no topo do ficheiro com a diretiva \texttt{@page}, por exemplo \verb|@page "/pathway/{Id:int}"|. As restrições de tipo (\texttt{:int}) fazem com que \texttt{/pathway/abc} não corresponda à página. Uma página pode declarar várias rotas: o construtor de percursos responde a três, uma por modo.

O \emph{Router} está envolvido num \texttt{AuthorizeRouteView}. Se a página tiver um atributo \texttt{[Authorize(Roles = \ldots)]} e o utilizador não tiver esse perfil, a aplicação redireciona-o para a página de entrada. O perfil é lido do \gls{jwt} guardado no navegador, pelo \texttt{CustomAuthStateProvider}.

Como o servidor devolve \texttt{index.html} a qualquer rota desconhecida, é possível abrir diretamente uma rota do cliente, por exemplo a partir de uma notificação ou de um marcador.

\subsection{Rotas da API}

Todos os controladores usam o mesmo modelo de rota,

\begin{center}
\verb|[Route("api/[controller]")]|,
\end{center}

\noindent em que \texttt{[controller]} é substituído pelo nome da classe sem o sufixo \texttt{Controller}. Cada ação acrescenta o seu segmento com o atributo do método HTTP. Por exemplo:

\begin{lstlisting}[language={[Sharp]C}, basicstyle=\footnotesize\ttfamily, caption={Rota de uma ação em \texttt{PathwayController}.}, label={lst:rota}]
[Route("api/[controller]")]        // -> api/pathway
public class PathwayController : ControllerBase
{
    [HttpGet("supervisor/{supervisorId}")]
    // GET api/pathway/supervisor/5
    public async Task<...> GetSupervisorPathways(int supervisorId)
}
\end{lstlisting}

As regras que se aplicam a toda a \gls{api} são:

\begin{itemize}
    \item As rotas não distinguem maiúsculas de minúsculas: \texttt{/api/Pathway} e \texttt{/api/pathway} são a mesma.
    \item Os segmentos entre chavetas (\texttt{\{id\}}) passam para o parâmetro com o mesmo nome; os corpos dos pedidos chegam em JSON (\texttt{[FromBody]} ou DTO).
    \item As respostas usam os códigos habituais: 200 com JSON, 204 sem conteúdo, 400 para pedidos inválidos, 401 sem \gls{jwt} válido e 404 quando o recurso não existe.
    \item A serialização ignora referências circulares (\texttt{ReferenceHandler.IgnoreCycles}), porque as entidades se referem umas às outras.
\end{itemize}

\section{Autenticação e autorização}

A autenticação usa \gls{jwt} assinados pelo próprio servidor:

\begin{enumerate}
    \item \texttt{POST /api/auth/login} verifica o email e o hash BCrypt da palavra-passe e devolve um token com o identificador, o nome, o email e o perfil do utilizador, válido durante 7 dias.
    \item O cliente guarda o token na \emph{session storage} do navegador. Por isso, a sessão termina quando o separador é fechado.
    \item O \texttt{AuthTokenHandler} junta o cabeçalho \texttt{Authorization: Bearer} a todos os pedidos feitos pelo \texttt{HttpClient} da aplicação.
    \item O servidor valida a assinatura, o emissor, o destinatário e a validade de cada token.
\end{enumerate}

A Figura~\ref{fig:sequencia} mostra este mecanismo num pedido real: gravar uma avaliação.

\begin{figure}[H]
    \centering
    \resizebox{\linewidth}{!}{% Diagrama de sequência: o supervisor grava uma avaliação.
\begin{tikzpicture}[
    ator/.style={draw=gInk, fill=gCard, rounded corners=3pt, font=\scriptsize\bfseries, align=center,
      text width=2.5cm, minimum height=0.9cm, inner sep=2pt, line width=0.6pt},
    vida/.style={draw=gInkSoft, dashed, line width=0.5pt},
    msg/.style={->, >={Stealth[length=1.8mm]}, draw=gInk, line width=0.6pt},
    ret/.style={->, >={Stealth[length=1.8mm, open]}, draw=gInkSoft, dashed, line width=0.5pt},
    txt/.style={font=\tiny, text=gInk, fill=white, inner sep=1pt, align=center},
    ativ/.style={draw=gInk, fill=gBrand!20, line width=0.5pt},
  ]
  % Participantes (x)
  \def\xA{0}\def\xB{2.8}\def\xC{5.6}\def\xD{8.4}\def\xE{11.2}\def\xF{14.0}
  \node[ator] (A) at (\xA,0.5) {Ficha de avaliação\\{\tiny\mdseries (Blazor)}};
  \node[ator] (B) at (\xB,0.5) {AuthToken\-Handler};
  \node[ator] (C) at (\xC,0.5) {Pipeline ASP.NET\\{\tiny\mdseries (JWT)}};
  \node[ator] (D) at (\xD,0.5) {Evaluation\-Controller};
  \node[ator] (E) at (\xE,0.5) {BadgeService};
  \node[ator] (F) at (\xF,0.5) {AppDbContext\\{\tiny\mdseries (SQL Server)}};
  \foreach \x in {\xA,\xB,\xC,\xD,\xE,\xF} \draw[vida] (\x,-0.1) -- (\x,-10.4);

  % Barras de ativação
  \fill[ativ] (\xD-0.08,-2.4) rectangle (\xD+0.08,-9.6);
  \draw[ativ] (\xD-0.08,-2.4) rectangle (\xD+0.08,-9.6);
  \fill[ativ] (\xE-0.08,-6.1) rectangle (\xE+0.08,-8.1);
  \draw[ativ] (\xE-0.08,-6.1) rectangle (\xE+0.08,-8.1);

  % Mensagens
  \draw[msg] (\xA,-1.0) -- node[txt, above=1pt] {POST /api/evaluation\\{[}lista de classificações{]}} (\xB,-1.0);
  \draw[msg] (\xB,-1.7) -- node[txt, above=1pt] {+ Authorization:\\Bearer \textless JWT\textgreater} (\xC,-1.7);
  \draw[msg] (\xC,-2.4) -- node[txt, above=1pt] {token válido:\\SaveEvaluations()} (\xD,-2.4);
  \draw[msg] (\xD,-3.2) -- node[txt, above=1pt, pos=0.75] {inserir ou atualizar ComponentEvaluations} (\xF,-3.2);
  \draw[msg] (\xD,-4.0) -- node[txt, above=1pt, pos=0.75] {ler módulos e componentes da inscrição} (\xF,-4.0);
  \draw[ret] (\xF,-4.5) -- (\xD,-4.5);
  % Auto-mensagem: calcular progresso
  \draw[msg] (\xD+0.08,-4.85) -- ++(0.7,0) -- ++(0,-0.4) -- (\xD+0.08,-5.25);
  \node[txt, anchor=west, text width=3.2cm, align=left] at (\xD+0.9,-5.05) {calcular o progresso do percurso e dos módulos};
  \draw[msg] (\xD,-6.1) -- node[txt, above=1pt] {EvaluateModuleBadgeAsync\\(por módulo alterado)} (\xE,-6.1);
  \draw[msg] (\xE,-6.8) -- node[txt, above=1pt] {UserBadge + Notification\\(se 65\% e ainda não tem)} (\xF,-6.8);
  \draw[msg] (\xD,-7.5) -- node[txt, above=1pt] {EvaluatePathwayBadgeAsync} (\xE,-7.5);
  \draw[ret] (\xE,-8.1) -- (\xD,-8.1);
  \draw[msg] (\xD,-8.8) -- node[txt, above=1pt] {Notification «Atualização de Avaliação»; SaveChanges} (\xF,-8.8);
  \draw[ret] (\xD,-9.6) -- node[txt, above=1pt] {200 OK} (\xA,-9.6);
\end{tikzpicture}
}
    \caption{Sequência de pedidos ao gravar uma avaliação.}
    \label{fig:sequencia}
\end{figure}

Todos os controladores exigem um utilizador autenticado (\texttt{[Authorize]}), com três exceções: o registo, o login e a transferência de ficheiros. A verificação do perfil, porém, só existe em dois endpoints. Nos restantes, a separação entre perfis é feita apenas pelas páginas do cliente, e os identificadores de utilizador chegam no URL em vez de serem lidos do token. Cumprir o requisito RNF12 exige verificar o perfil e a posse dos dados no servidor, em cada pedido.

\section{Catálogo de endpoints}
\label{sec:endpoints}

A \gls{api} tem 54 endpoints, distribuídos por 14 controladores (Tabela~\ref{tab:controladores}). A Tabela~\ref{tab:endpoints} lista-os todos. Na coluna \emph{Acesso}, \emph{Autenticado} significa qualquer utilizador com um \gls{jwt} válido.

\begin{table}[H]
    \centering
    \small
    \begin{tabularx}{\linewidth}{lrX}
        \toprule
        \textbf{Controlador} & \textbf{Endpoints} & \textbf{Responsabilidade} \\
        \midrule
        \texttt{Auth} & 4 & Registo, login e palavras-passe. \\
        \texttt{User} & 2 & Listagem de utilizadores. \\
        \texttt{PathwayTemplate} & 4 & Moldes de percurso. \\
        \texttt{ModuleTemplate} & 4 & Módulos dos moldes. \\
        \texttt{ComponentTemplate} & 4 & Componentes dos moldes. \\
        \texttt{BadgeTemplate} & 4 & Badges dos moldes. \\
        \texttt{Pathway} & 7 & Percursos e cópia a partir de moldes. \\
        \texttt{Module} & 6 & Módulos e datas das fases. \\
        \texttt{ModuleComponent} & 4 & Componentes dos módulos. \\
        \texttt{Enrollment} & 7 & Inscrições e cópia do percurso por aluno. \\
        \texttt{Evaluation} & 3 & Avaliações, progresso e relatório \gls{pdf}. \\
        \texttt{Badges} & 1 & Badges obtidas por um supervisionado. \\
        \texttt{Notification} & 2 & Notificações. \\
        \texttt{Files} & 2 & Carregamento e transferência de documentos. \\
        \midrule
        \textbf{Total} & \textbf{54} & \\
        \bottomrule
    \end{tabularx}
    \caption{Controladores da API.}
    \label{tab:controladores}
\end{table}

\newcommand{\ep}[4]{\texttt{#1} & \texttt{#2} & #3 & #4 \\}
\newcommand{\grupo}[1]{\midrule\multicolumn{4}{l}{\textbf{#1}} \\*}

\begingroup\footnotesize
\begin{longtable}{l >{\raggedright\arraybackslash}p{0.40\linewidth} >{\raggedright\arraybackslash}p{0.27\linewidth} >{\raggedright\arraybackslash}p{0.12\linewidth}}
\toprule
\textbf{Método} & \textbf{Rota} & \textbf{Descrição} & \textbf{Acesso} \\
\midrule
\endfirsthead
\toprule
\textbf{Método} & \textbf{Rota} & \textbf{Descrição} & \textbf{Acesso} \\
\midrule
\endhead
\bottomrule
\endfoot
\bottomrule
\caption{Catálogo de endpoints da API.}\label{tab:endpoints}\\
\endlastfoot
\multicolumn{4}{l}{\textbf{AuthController} --- \texttt{/api/auth}} \\*
\ep{POST}{/\allowbreak{}api/\allowbreak{}auth/\allowbreak{}register}{Criar utilizador}{Público}
\ep{POST}{/\allowbreak{}api/\allowbreak{}auth/\allowbreak{}login}{Autenticar e obter o \gls{jwt}}{Público}
\ep{POST}{/\allowbreak{}api/\allowbreak{}auth/\allowbreak{}change-password}{Alterar a própria palavra-passe}{Autenticado}
\ep{POST}{/\allowbreak{}api/\allowbreak{}auth/\allowbreak{}reset-password}{Repor a palavra-passe de um utilizador}{Admin}
\grupo{UserController --- \texttt{/api/user}}
\ep{GET}{/\allowbreak{}api/\allowbreak{}user}{Utilizadores que não são administradores}{Autenticado}
\ep{GET}{/\allowbreak{}api/\allowbreak{}user/\allowbreak{}students}{Supervisionados}{Autenticado}
\grupo{PathwayTemplateController --- \texttt{/api/pathwaytemplate}}
\ep{GET}{/\allowbreak{}api/\allowbreak{}pathwaytemplate}{Listar moldes}{Autenticado}
\ep{GET}{/\allowbreak{}api/\allowbreak{}pathwaytemplate/\allowbreak{}\{id\}}{Obter um molde}{Autenticado}
\ep{POST}{/\allowbreak{}api/\allowbreak{}pathwaytemplate}{Criar molde}{Autenticado}
\ep{DELETE}{/\allowbreak{}api/\allowbreak{}pathwaytemplate/\allowbreak{}\{id\}}{Apagar molde}{Autenticado}
\grupo{ModuleTemplateController --- \texttt{/api/moduletemplate}}
\ep{GET}{/\allowbreak{}api/\allowbreak{}moduletemplate/\allowbreak{}\{id\}}{Obter módulo do molde}{Autenticado}
\ep{POST}{/\allowbreak{}api/\allowbreak{}moduletemplate}{Criar módulo do molde}{Autenticado}
\ep{PUT}{/\allowbreak{}api/\allowbreak{}moduletemplate/\allowbreak{}\{id\}}{Atualizar módulo do molde}{Autenticado}
\ep{DELETE}{/\allowbreak{}api/\allowbreak{}moduletemplate/\allowbreak{}\{id\}}{Apagar módulo do molde}{Autenticado}
\grupo{ComponentTemplateController --- \texttt{/api/componenttemplate}}
\ep{GET}{/\allowbreak{}api/\allowbreak{}componenttemplate/\allowbreak{}\{id\}}{Obter componente do molde}{Autenticado}
\ep{POST}{/\allowbreak{}api/\allowbreak{}componenttemplate}{Criar componente do molde}{Autenticado}
\ep{PUT}{/\allowbreak{}api/\allowbreak{}componenttemplate/\allowbreak{}\{id\}}{Atualizar componente do molde}{Autenticado}
\ep{DELETE}{/\allowbreak{}api/\allowbreak{}componenttemplate/\allowbreak{}\{id\}}{Apagar componente do molde}{Autenticado}
\grupo{BadgeTemplateController --- \texttt{/api/badgetemplate}}
\ep{GET}{/\allowbreak{}api/\allowbreak{}badgetemplate/\allowbreak{}pathway/\allowbreak{}\{pathwayTemplateId\}}{Badges de um molde}{Autenticado}
\ep{POST}{/\allowbreak{}api/\allowbreak{}badgetemplate}{Criar badge do molde}{Autenticado}
\ep{PUT}{/\allowbreak{}api/\allowbreak{}badgetemplate/\allowbreak{}\{id\}}{Atualizar badge do molde}{Autenticado}
\ep{DELETE}{/\allowbreak{}api/\allowbreak{}badgetemplate/\allowbreak{}\{id\}}{Apagar badge do molde}{Autenticado}
\grupo{PathwayController --- \texttt{/api/pathway}}
\ep{POST}{/\allowbreak{}api/\allowbreak{}pathway}{Criar percurso (copia o molde, se indicado)}{Autenticado}
\ep{GET}{/\allowbreak{}api/\allowbreak{}pathway/\allowbreak{}\{id\}}{Obter percurso}{Autenticado}
\ep{GET}{/\allowbreak{}api/\allowbreak{}pathway/\allowbreak{}supervisor/\allowbreak{}\{supervisorId\}}{Percursos de um supervisor}{Autenticado}
\ep{GET}{/\allowbreak{}api/\allowbreak{}pathway/\allowbreak{}\{pathwayId\}/enrollments}{Inscrições de um percurso}{Autenticado}
\ep{GET}{/\allowbreak{}api/\allowbreak{}pathway/\allowbreak{}builder/\allowbreak{}\{id\}}{Cópia de um aluno, para o construtor}{Autenticado}
\ep{PUT}{/\allowbreak{}api/\allowbreak{}pathway/\allowbreak{}\{id\}/title}{Mudar o título}{Autenticado}
\ep{DELETE}{/\allowbreak{}api/\allowbreak{}pathway/\allowbreak{}\{id\}}{Apagar percurso}{Autenticado}
\grupo{ModuleController --- \texttt{/api/module}}
\ep{POST}{/\allowbreak{}api/\allowbreak{}module}{Criar módulo}{Autenticado}
\ep{GET}{/\allowbreak{}api/\allowbreak{}module/\allowbreak{}\{id\}}{Obter módulo}{Autenticado}
\ep{GET}{/\allowbreak{}api/\allowbreak{}module/\allowbreak{}\{moduleId\}/student/\{studentId\}}{Módulo com as avaliações de um aluno}{Autenticado}
\ep{PUT}{/\allowbreak{}api/\allowbreak{}module/\allowbreak{}\{id\}}{Atualizar módulo}{Autenticado}
\ep{PUT}{/\allowbreak{}api/\allowbreak{}module/\allowbreak{}\{id\}/timeline}{Atualizar as datas das fases de um aluno}{Autenticado}
\ep{DELETE}{/\allowbreak{}api/\allowbreak{}module/\allowbreak{}\{id\}}{Apagar módulo}{Autenticado}
\grupo{ModuleComponentController --- \texttt{/api/modulecomponent}}
\ep{POST}{/\allowbreak{}api/\allowbreak{}modulecomponent}{Criar componente}{Autenticado}
\ep{GET}{/\allowbreak{}api/\allowbreak{}modulecomponent/\allowbreak{}\{id\}}{Obter componente}{Autenticado}
\ep{PUT}{/\allowbreak{}api/\allowbreak{}modulecomponent/\allowbreak{}\{id\}}{Atualizar componente}{Autenticado}
\ep{DELETE}{/\allowbreak{}api/\allowbreak{}modulecomponent/\allowbreak{}\{id\}}{Apagar componente}{Autenticado}
\grupo{EnrollmentController --- \texttt{/api/enrollment}}
\ep{POST}{/\allowbreak{}api/\allowbreak{}enrollment}{Inscrever aluno (copia o percurso)}{Autenticado}
\ep{GET}{/\allowbreak{}api/\allowbreak{}enrollment/\allowbreak{}supervisor/\allowbreak{}\{supervisorId\}}{Inscrições dos percursos de um supervisor}{Autenticado}
\ep{GET}{/\allowbreak{}api/\allowbreak{}enrollment/\allowbreak{}student/\allowbreak{}\{studentId\}}{Inscrições de um aluno}{Autenticado}
\ep{GET}{/\allowbreak{}api/\allowbreak{}enrollment/\allowbreak{}student/\allowbreak{}\{studentId\}/pathway/\{pathwayId\}}{Detalhe de uma inscrição}{Autenticado}
\ep{GET}{/\allowbreak{}api/\allowbreak{}enrollment/\allowbreak{}management}{Alunos, para a gestão de utilizadores}{Supervisor, Admin}
\ep{PUT}{/\allowbreak{}api/\allowbreak{}enrollment/\allowbreak{}\{id\}/star}{Marcar ou desmarcar como favorita}{Autenticado}
\ep{PUT}{/\allowbreak{}api/\allowbreak{}enrollment/\allowbreak{}\{id\}/hidden}{Ocultar ou mostrar no painel}{Autenticado}
\grupo{EvaluationController --- \texttt{/api/evaluation}}
\ep{GET}{/\allowbreak{}api/\allowbreak{}evaluation/\allowbreak{}enrollment/\allowbreak{}\{id\}}{Inscrição com as avaliações}{Autenticado}
\ep{POST}{/\allowbreak{}api/\allowbreak{}evaluation}{Gravar avaliações; recalcular progresso e badges}{Autenticado}
\ep{GET}{/\allowbreak{}api/\allowbreak{}evaluation/\allowbreak{}student/\allowbreak{}\{studentId\}/pdf}{Relatório global em \gls{pdf}}{Autenticado}
\grupo{BadgesController --- \texttt{/api/badges}}
\ep{GET}{/\allowbreak{}api/\allowbreak{}badges/\allowbreak{}student/\allowbreak{}\{studentId\}}{Badges de um aluno}{Autenticado}
\grupo{NotificationController --- \texttt{/api/notification}}
\ep{GET}{/\allowbreak{}api/\allowbreak{}notification/\allowbreak{}user/\allowbreak{}\{userId\}}{Notificações de um utilizador}{Autenticado}
\ep{PUT}{/\allowbreak{}api/\allowbreak{}notification/\allowbreak{}\{id\}/read}{Marcar como lida}{Autenticado}
\grupo{FilesController --- \texttt{/api/files}}
\ep{POST}{/\allowbreak{}api/\allowbreak{}files/\allowbreak{}upload}{Carregar um documento}{Autenticado}
\ep{GET}{/\allowbreak{}api/\allowbreak{}files/\allowbreak{}download/\allowbreak{}\{fileName\}}{Transferir um documento}{Público}
\end{longtable}
\endgroup

\section{Dados}

A base de dados é gerida com Entity Framework Core em modo \emph{code-first}: as tabelas são geradas a partir das classes da biblioteca partilhada e evoluem através de migrações. São 15 tabelas, organizadas em três camadas que correspondem às duas cópias descritas na Secção~\ref{sec:estilo}: moldes, percursos partilhados e cópias por aluno. As chaves estrangeiras entre inscrições, avaliações e utilizadores usam \texttt{Restrict}, para que apagar um utilizador ou um módulo não apague em cascata o histórico de avaliações.

Os documentos \gls{pdf} carregados são guardados fora da pasta pública do servidor, em \texttt{SecureUploads/}, com um nome aleatório (GUID), e só são acessíveis através do endpoint de transferência.

\porcompletar{diagrama entidade-relação das 15 tabelas, se o orientador o pedir.}
